\chapter{Traffic engineering methodology}
\label{chap:methodology}

This chapter identifies the source, definition and implementation status of the principal PySig calculations. PySig combines established traffic-engineering relationships with original deterministic simulation and optimisation algorithms. The distinction is material: a familiar engineering quantity such as degree of saturation may be calculated analytically, while a queue or accessible service capacity may be generated by the PySig time-sliced kernel.

\section{Method provenance}
\begin{longtable}{p{0.31\textwidth}p{0.60\textwidth}}
\toprule
\textbf{PySig quantity} & \textbf{Method and provenance}\\
\midrule\endhead
Saturation flow & Direct measured input or the implemented RR67-style unopposed saturation-flow relationship. Measured critical-lane values are preferred where suitable observations are available.\\
Signal green & Sum of phase durations during which the signal group is active. Intergreen is excluded.\\
Effective green & Signal green minus start displacement plus end displacement, bounded to zero and the cycle time.\\
Nominal capacity & Conventional fixed-time relationship $c=s(g/C)$.\\
Nominal degree of saturation & Conventional relationship $x=q/c$.\\
Physical lane queue & PySig deterministic, time-sliced conservation model.\\
Shared-lane interaction & PySig FIFO-composition approximation for movements sharing one physical lane.\\
Delay & Time integral, or discrete area, under the simulated lane-queue curve.\\
Queue metres & Simulated queue in PCU multiplied by configured PCU length.\\
Route choice & Multinomial-logit target allocation.\\
Assignment update & Method of successive averages.\\
OD conservation & Exact accounting constraint for every origin-destination pair.\\
Signal optimisation & Bounded candidate enumeration and sequential controller optimisation.\\
Network propagation & PySig time-sliced link travel, storage and simplified dispersion approximation.\\
\bottomrule
\end{longtable}

\section{Signal green, intergreen and displacement}
For signal group $j$, signal green $G_j$ is the sum of the durations of phases in which the group is active:
\begin{equation}
G_j=\sum_{p:\,j\in p}G_p.
\end{equation}
Intergreen is part of the controller cycle but is not included in signal green:
\begin{equation}
C=\sum_p(G_p+I_p).
\end{equation}

PySig defines start displacement $d_s$ as effective service lost after displayed green begins. End displacement $d_e$ is effective service gained after displayed green ends. Effective green is therefore
\begin{equation}
g_j=G_j-d_s+d_e,
\end{equation}
with the calculated result bounded to $0\leq g_j\leq C$.

For a 60-second cycle with 55 seconds of signal green, 5 seconds of intergreen, 3 seconds of start displacement and 3 seconds of end displacement:
\begin{align}
g &= 55-3+3\\
  &= 55\ \text{s}.
\end{align}
Equal start and end displacements cancel. An asymmetric case, such as 3 seconds at the start and 2 seconds at the end, gives 54 seconds of effective green.

\section{Nominal lane capacity}
Nominal fixed-time capacity is
\begin{equation}
c_n=s\frac{g}{C},
\end{equation}
where $s$ is saturation flow in PCU/h of effective green, $g$ is effective green and $C$ is cycle time.

For $s=1900$ PCU/h, $g=55$ s and $C=60$ s:
\begin{align}
c_n &=1900\frac{55}{60}\\
    &=1741.67\ \text{PCU/h}.
\end{align}

\section{Nominal degree of saturation}
The conventional nominal degree of saturation is
\begin{equation}
x_n=\frac{q}{c_n}=\frac{qC}{sg}.
\end{equation}
For demand of 1000 PCU/h in the preceding example:
\begin{equation}
x_n=\frac{1000}{1741.67}=0.57.
\end{equation}

\section{Simulated service capacity and simulated saturation}
PySig separately reports service made accessible by the time-sliced model. Shared-lane blocking, finite downstream storage and network spillback may make this lower than nominal capacity. The corresponding simulated degree of saturation is
\begin{equation}
x_s=\frac{q}{c_s}.
\end{equation}
Nominal and simulated quantities answer different questions. Nominal capacity is independently checkable from signal timing and saturation flow. Simulated service capacity reports what the current PySig operational kernel made available.

\section{Lane queue conservation}
For each physical lane and time step $\Delta t$:
\begin{equation}
Q_{t+\Delta t}=Q_t+A_t-D_t,
\end{equation}
subject to finite storage
\begin{equation}
0\leq Q_t\leq K.
\end{equation}
Movements sharing a \code{from_lane_id} contribute to one physical queue:
\begin{equation}
Q_{\mathrm{lane}}=\sum_{m\in M_{\mathrm{lane}}}Q_m.
\end{equation}
Saturation flow and storage are applied once per physical lane, not once per movement.

\section{Delay}
Lane delay is the area under the queue-length curve. In the time-sliced implementation:
\begin{equation}
D\approx\sum_t Q_t\Delta t.
\end{equation}
PySig reports total delay in PCU-hours and mean delay in seconds per departed PCU. Movement delay is an allocation of the physical lane delay to movement subqueue composition; it is not a second independent delay process.

\section{Final queue, unserved demand and overflow}
These values should not be conflated:
\begin{itemize}
\item \textbf{Final queue} is traffic remaining physically queued at the end of the modelled horizon.
\item \textbf{Unserved demand} is accepted traffic that has not departed by the end of the horizon.
\item \textbf{Storage overflow} is arriving traffic that could not enter the modelled lane because physical storage was full.
\end{itemize}

\section{Output precision}
PySig retains full floating-point precision internally. Operational outputs are rounded to two decimal places only during JSON and CSV serialisation. Convergence and conservation diagnostics retain six decimal places so that output formatting does not conceal numerical residuals.
